\chapter{Fungal Culture}

\section*{What Is This Test?}
Fungal culture (mycology culture) is the microbiologic method for isolating, identifying, and determining antifungal susceptibility of fungi from clinical specimens. Fungi are eukaryotic organisms with cell walls containing chitin, glucans, and mannoproteins. Clinically significant fungi are broadly categorized as yeasts (unicellular, reproduce by budding: \textit{Candida} species, \textit{Cryptococcus neoformans/gattii}), molds (multicellular, filamentous, produce hyphae and spores: \textit{Aspergillus} species, Mucorales/Zygomycetes, \textit{Fusarium}, \textit{Scedosporium}, dermatophytes), and dimorphic fungi (exist as mold at 25--30\textdegree{}C in the environment and as yeast/spherule at 37\textdegree{}C in human tissue: \textit{Histoplasma capsulatum}, \textit{Coccidioides immitis/posadasii}, \textit{Blastomyces dermatitidis/gilchristii}, \textit{Paracoccidioides brasiliensis}, \textit{Sporothrix schenckii}, \textit{Talaromyces} [formerly \textit{Penicillium}] \textit{marneffei}). Invasive fungal infections (IFIs) are life-threatening and occur predominantly in immunocompromised patients: HSCT and SOT recipients, patients with hematologic malignancies and prolonged neutropenia, those on high-dose corticosteroids or biologic immunosuppressants, and premature neonates. Fungal culture sensitivity varies by organism and specimen type, typically ranging from 30--80\%, and growth is often slow (days to weeks for molds and dimorphic fungi). Non-culture methods (fungal biomarkers: serum beta-D-glucan, galactomannan, \textit{Cryptococcus} antigen; PCR) complement culture for early diagnosis of IFIs.

\section*{Alternative Names}
Mycology Culture; Mold Culture; Yeast Culture; Fungal Isolation; Fungal C\&S; Dermatophyte Culture; Dimorphic Fungus Culture

\section*{Principle}
Specimens are inoculated onto a battery of fungal media: Sabouraud dextrose agar (SDA, with gentamicin and chloramphenicol to inhibit bacteria; the standard general-purpose fungal medium), brain-heart infusion (BHI) agar with 5--10\% sheep blood (for dimorphic fungi recovery at 25--30\textdegree{}C), inhibitory mold agar (IMA, with gentamicin and chloramphenicol, for general mold isolation), dermatophyte test medium (DTM; contains phenol red pH indicator that turns red when dermatophytes grow, providing rapid presumptive identification), and Mycosel/Mycobiotic agar (SDA with cycloheximide to inhibit saprophytic molds and chloramphenicol; useful for isolating dermatophytes and dimorphic fungi, but cycloheximide inhibits \textit{Aspergillus}, Mucorales, \textit{Cryptococcus}, and many opportunistic molds, so both non-selective and selective media must be inoculated). Incubation: Yeasts at 30--35\textdegree{}C for 24--72 hours. Molds and dimorphic fungi: 25--30\textdegree{}C for up to 4--6 weeks (dimorphic fungi may require 2--6 weeks for initial growth). The culture is examined daily for the first week and weekly thereafter. Yeasts: Cream-colored, pasty colonies. Molds: Filamentous, wooly, cottony, or powdery colonies with various colors (green, black, white, brown, yellow, pink). Species identification for yeasts: MALDI-TOF mass spectrometry (rapid, accurate), germ tube test (\textit{C. albicans} produces germ tubes in 2--3 hours at 37\textdegree{}C in serum---a rapid presumptive test), chromogenic agar (CHROMagar Candida --- \textit{C. albicans}: green; \textit{C. tropicalis}: blue; \textit{C. krusei}: pink/rough; \textit{C. glabrata}: purple), carbohydrate assimilation (API 20C AUX, VITEK 2 YST) and urease production. Species identification for molds: Macroscopic (colony morphology, growth rate, color, texture) and microscopic (lactophenol cotton blue tease mount or slide culture: hyphal morphology, septation, pigmentation, and reproductive structures---conidia, sporangiospores, phialides, metulae, vesicles, etc.). Identification of dimorphic fungi involves confirming conversion from the mold form (25\textdegree{}C) to the yeast or spherule form (37\textdegree{}C) on enriched media, or by DNA probe (Gen-Probe AccuProbe) or MALDI-TOF. Definitive identification often requires sequencing of the internal transcribed spacer (ITS) region of ribosomal DNA (ITS1-5.8S-ITS2) and/or the D1/D2 region of 28S rDNA, which is the molecular gold standard. Antifungal susceptibility testing (AFST): CLSI M27 (yeasts) and M38 (molds) broth microdilution methods determine minimum inhibitory concentrations (MICs) for azoles (fluconazole, voriconazole, posaconazole, isavuconazole), echinocandins (caspofungin, micafungin, anidulafungin), amphotericin B, and flucytosine. AFST is indicated for: invasive isolates from sterile sites, isolates from patients with refractory or breakthrough infections, \textit{Candida} species with known or potential resistance (\textit{C. glabrata}, \textit{C. krusei}, \textit{C. auris}), and all \textit{Aspergillus} isolates when azole therapy is planned (due to emerging azole resistance through \textit{cyp51A} mutations).

\section*{History}
Fungal culture was pioneered by Raymond Sabouraud in the 1890s for the study of dermatophytes. The dimorphic nature of \textit{Histoplasma} was recognized by Samuel Darling in 1906 and \textit{Coccidioides} by Emmet Rixford and T.C. Gilchrist in 1896. Amphotericin B, the first systemic antifungal, was introduced in the 1950s. Fluconazole and itraconazole (triazoles), introduced in the 1990s, expanded therapeutic options. The echinocandins (caspofungin 2001, micafungin 2005, anidulafungin 2006) provided a new class targeting the fungal cell wall. MALDI-TOF MS for yeast identification became standard in the 2010s, dramatically reducing time to identification from 48--72 hours to minutes. \textit{Candida auris}, an emerging multidrug-resistant yeast, was first described in 2009 and has since caused nosocomial outbreaks globally, highlighting the importance of species-level identification and AFST. The WHO Fungal Priority Pathogens List (2022) identified \textit{C. auris}, \textit{Aspergillus fumigatus}, \textit{Cryptococcus neoformans}, and \textit{Candida albicans} as critical priority pathogens.

\section*{Why This Test Is Ordered}
\begin{itemize}
  \item Diagnosis of invasive candidiasis (candidemia, intra-abdominal candidiasis, hepatosplenic candidiasis) in at-risk patients: ICU patients with prolonged stay, central venous catheters, broad-spectrum antibiotics, total parenteral nutrition, abdominal surgery, and immunocompromised status \cite{pappas2016clinical}. Blood cultures are positive in only 50--70\% of cases of invasive candidiasis
  \item Diagnosis of invasive aspergillosis (pulmonary, sinus, CNS, disseminated) in severely immunocompromised patients: prolonged neutropenia (>10--14 days), HSCT (particularly with GVHD), SOT (lung > heart > liver > kidney), and chronic granulomatous disease. BAL culture plus galactomannan (BAL and serum) is the standard diagnostic approach; culture sensitivity is 30--60\% \cite{patterson2016practice}
  \item Diagnosis of mucormycosis (zygomycosis) caused by Mucorales (\textit{Rhizopus}, \textit{Mucor}, \textit{Lichtheimia/Absidia}, \textit{Cunninghamella}) in immunocompromised patients, particularly those with diabetic ketoacidosis, iron overload (deferoxamine therapy), or HSCT. The organism is angioinvasive, causing tissue necrosis, infarction, and rapid progression \cite{cornely2019global}. Culture sensitivity is disappointingly low (25--50\%), and histopathology of tissue biopsy is often diagnostic (broad, ribbon-like, pauci-septate hyphae with right-angle branching invading blood vessels)
  \item Diagnosis of endemic dimorphic mycoses (\textit{Histoplasma}, \textit{Coccidioides}, \textit{Blastomyces}) in patients with travel to or residence in endemic areas presenting with pneumonia, mediastinal lymphadenopathy, or disseminated disease. Culture is the gold standard but slow (2--6 weeks); antigen detection (urine/serum) and serology provide earlier diagnosis
  \item Diagnosis of cryptococcal meningitis (CSF culture, India ink, and cryptococcal antigen/CrAg) in HIV/AIDS patients with CD4 <100 cells/$\mu$L and in SOT recipients
  \item Diagnosis of dermatophyte infections (tinea corporis, tinea cruris, tinea pedis, tinea capitis, onychomycosis) when the diagnosis is uncertain or treatment is refractory
  \item Antifungal susceptibility testing to guide therapy in drug-resistant fungi, breakthrough infections, and epidemiological surveillance
\end{itemize}

\section*{Primary vs Secondary Test}
Fungal culture is the primary test for diagnosis of fungal infections. It is the gold standard for: (1) definitive species identification of fungi; (2) antifungal susceptibility testing; (3) recovery of viable organisms for epidemiologic studies. However, fungal culture has limited sensitivity (30--80\%, depending on the organism and specimen type) and slow turnaround (days to weeks). Non-culture biomarkers (galactomannan, beta-D-glucan, cryptococcal antigen, \textit{Histoplasma} antigen, \textit{Blastomyces} antigen) complement culture and are often positive earlier.

\section*{How This Test Is Performed}
Blood cultures for fungi: Inoculated into aerobic blood culture bottles or dedicated fungal blood culture bottles (Isolator tube, BACTEC Myco/F Lytic). Standard aerobic bottles detect most \textit{Candida} species. Fungal-specific bottles improve recovery of \textit{Histoplasma}, \textit{Cryptococcus}, and \textit{Malassezia} (which requires lipid supplementation). Incubation for 5--7 days is standard for yeasts; extended incubation (2--4 weeks) for molds and dimorphic fungi may be requested. Respiratory specimens: BAL, bronchial wash, sputum, endotracheal aspirate. Processed similarly to bacterial culture but plated on fungal media. CSF: Large volumes (5--10 mL) improve sensitivity. Processed by centrifugation and inoculation onto fungal media. India ink preparation provides rapid presumptive identification of encapsulated \textit{Cryptococcus}. Tissue biopsy: Homogenized tissue is plated; histopathology (Gomori methenamine silver/GMS, periodic acid-Schiff/PAS, hematoxylin and eosin) is performed on adjacent tissue sections. Skin, hair, and nails for dermatophytes: Placed on DTM and SDA with cycloheximide and chloramphenicol; incubated at 25--30\textdegree{}C for up to 4 weeks. All fungal cultures should be processed in a biosafety cabinet (BSC Class II) because dimorphic fungi in the mold form (\textit{Histoplasma}, \textit{Coccidioides}, \textit{Blastomyces}) are highly infectious by the airborne route and are classified as Biosafety Level 3 (BSL-3) agents in their mold form.

\section*{What Preparation Is Needed}
No special patient preparation beyond that for the specific specimen collection (blood draw, bronchoscopy, lumbar puncture, tissue biopsy). Skin scrapings for dermatophyte culture should be collected from the active border of the lesion after cleansing with 70\% alcohol (to reduce bacterial contamination).

\section*{Contraindications and Precautions}
Fungal culture is slow, and therapy decisions in invasive fungal infections must often be made empirically before culture results return (empiric/preemptive antifungal therapy based on host risk factors, clinical presentation, imaging, and fungal biomarkers) \cite{walsh2008treatment}. A negative fungal culture does NOT exclude invasive fungal infection; culture sensitivity for invasive aspergillosis is only 30--60\%, and for mucormycosis 25--50\%. Organism viability may be lost during specimen transport, particularly for Mucorales (the non-septate hyphae are fragile and easily damaged during tissue grinding/homogenization). Histopathologic examination of tissue (GMS stain) demonstrating fungal hyphae invading tissue is diagnostic even when culture is negative. Cycloheximide-containing media (Mycosel, Mycobiotic) inhibits many clinically significant molds (\textit{Aspergillus}, Mucorales, \textit{Cryptococcus}, \textit{Fusarium}, \textit{Scedosporium}); both cycloheximide-containing and non-selective media must be inoculated. \textit{Candida} species are colonizers of the respiratory and gastrointestinal tracts; isolation of \textit{Candida} from sputum, BAL, or stool does NOT establish the diagnosis of invasive candidiasis and should NOT be treated (treatment of airway colonization is a common antibiotic stewardship error). Isolation of \textit{Candida} from a normally sterile site (blood, CSF, peritoneal fluid, pleural fluid) always represents infection.

\section*{Risks}
No risks from the culture test itself. Risks are associated with specimen collection. Laboratory personnel must handle mold cultures (suspected dimorphic fungi) in a BSC with appropriate personal protective equipment, as \textit{Coccidioides immitis} and other dimorphic molds are highly infectious via aerosol.

\section*{What Happens After the Test}
Yeast identification: 24--72 hours after detection of growth. Mold identification: 3 days to 4 weeks (average 1--2 weeks). Antifungal susceptibility testing: 48--96 hours after isolation of a pure culture. Positive yeast from blood cultures: empiric antifungal therapy with an echinocandin (caspofungin, micafungin, anidulafungin) should be initiated. Species identification and susceptibility results guide de-escalation to fluconazole (for susceptible \textit{C. albicans}) or continuation/modification for resistant species (\textit{C. glabrata}, \textit{C. krusei}, \textit{C. auris}). All patients with candidemia require ophthalmologic examination for endophthalmitis and echocardiography (preferably TEE) for endocarditis.

\section*{Understanding Results}

\subsection*{Below Normal}
\begin{itemize}
  \item \textbf{No fungus isolated at 4 weeks:} No fungal growth on any media. Does not exclude fungal infection (culture sensitivity 30--80\% depending on organism and specimen). Biomarker tests (galactomannan, beta-D-glucan) and histopathology should be correlated. For dermatophyte cultures, incubation may be extended to 4--6 weeks
\end{itemize}

\subsection*{Above Normal}
\begin{itemize}
  \item \textbf{Yeast isolated:} \textit{Candida} species (most common), \textit{Cryptococcus neoformans/gattii}, other pathogenic yeasts. Species-level identification and antifungal susceptibility are reported. \textit{Candida albicans}: typically azole-susceptible. \textit{C. glabrata}: inherently reduced fluconazole susceptibility; echinocandin preferred. \textit{C. krusei}: inherently fluconazole-resistant. \textit{C. auris}: frequently multidrug-resistant (azole, amphotericin B, and/or echinocandin); urgent public health notification required
  \item \textbf{Mold isolated:} Species identification and, when indicated, antifungal susceptibility testing performed. Common molds: \textit{Aspergillus fumigatus} (most common cause of invasive aspergillosis), \textit{A. flavus}, \textit{A. niger}, \textit{A. terreus} (intrinsically resistant to amphotericin B). Mucorales group: \textit{Rhizopus, Mucor, Lichtheimia/Absidia, Cunninghamella}. Rapid growers vs. slow growers, thermotolerance (growth at 37\textdegree{}C vs. 42\textdegree{}C for \textit{A. fumigatus} vs. non-\textit{fumigatus}) assists in identification
  \item \textbf{Dimorphic fungus isolated:} Confirmed by conversion to the yeast/spherule phase at 37\textdegree{}C and/or DNA probe. \textit{Histoplasma capsulatum, Coccidioides immitis/posadasii}, and \textit{Blastomyces dermatitidis} are nationally notifiable in many jurisdictions. The laboratory must be notified when dimorphic fungi are suspected so that proper BSL-2+ (mold phase) precautions are observed
\end{itemize}

\section*{Normal Results}
\textbf{No fungus isolated} (at 4 weeks for routine cultures; up to 6 weeks for dermatophytes and some dimorphic fungi). Negative culture does not exclude fungal infection.

\section*{Follow-up Tests}
\begin{itemize}
  \item \textbf{Fungal biomarkers:} (1) Serum and BAL galactomannan (aspergillosis): index $\geq$0.5 (serum) or $\geq$1.0 (BAL) is positive. (2) Serum beta-D-glucan (candidiasis, aspergillosis, \textit{Pneumocystis}, fusariosis): $\geq$80 pg/mL is positive; sensitive but nonspecific (false positives with albumin, IVIG, hemodialysis membranes, gauze). (3) Cryptococcal antigen (CrAg): serum and CSF; highly sensitive and specific for cryptococcosis. (4) \textit{Histoplasma} urine antigen: >90\% sensitivity in disseminated histoplasmosis
  \item \textbf{Antifungal susceptibility testing (broth microdilution per CLSI M27/M38):} MICs for azoles, echinocandins, and amphotericin B are determined for isolates from invasive disease and in the setting of clinical failure or breakthrough infection \cite{clsi2017m27,clsi2017m38}
  \item \textbf{DNA sequencing (ITS + D1/D2 region):} Definitive species identification for molds when morphologic identification is inconclusive; the molecular gold standard
  \item \textbf{Histopathology (GMS, PAS stains):} Tissue biopsy with demonstration of fungal elements invading tissue is diagnostic, even when culture is negative. Histopathology provides rapid presumptive classification (yeast vs. mold vs. dimorphic yeast, hyphal morphology)
  \item \textbf{Blood cultures (repeat):} For candidemia, daily blood cultures until clearance is documented (negative cultures for $\geq$48 hours). Persistent candidemia beyond 3--5 days of appropriate therapy suggests an uncontrolled source, endocarditis, or antifungal resistance
\end{itemize}

\section*{References}
\bibliographystyle{unsrtnat}
\bibliography{references}

\clearpage
